\documentclass{article}
\usepackage{amsmath}
\usepackage{amssymb}
\usepackage{booktabs}
\usepackage{array}
\usepackage[a4paper,margin=1in]{geometry}

\newcolumntype{L}[1]{>{\raggedright\arraybackslash}p{#1}}

\title{Choosing the Promoter-Mining Ensemble\\ \large How six DNA models were compared, and why three were deployed}
\date{September 2026}
\author{HKUST iGEM 2026 Dry Lab}

\begin{document}
\maketitle

\section{Question and protocol}
We compared curated-window classification with restricted genomic ranking to choose the promoter-mining ensemble. Tables are historical records, not new evaluations. Curated AP depends on split, prevalence and weighting; the displayed six-model comparison mixes checkpoints and NTv3's 384-bp context.

Silver positives are heuristically labelled, putative promoter-positive upstream DNA. Labels come from the annotation and genomic position of predicted housekeeping loci, not measured promoter activity or mapped TSS. Putative-negative windows are not experimentally established promoter-free DNA.

The step-10 benchmark covers published TSS with same-contig/strand 81-bp windows expanded by 40 bp at each end. Unmatched starts merge at gaps of at most 81 bp. Recall counts covered catalogue TSS; precision proxy is covered TSS/(covered TSS+unmatched peaks), not ordinary matched-peak precision. Incomplete truth leaves unmatched calls unvalidated. Both benchmark strains supplied silver positives and putative-negative training windows; their published TSS labels were withheld.

\section{Curated classification}
Earlier T21\_final/plain-0.6 tests have 967/3,318 positives (no-skill AP about 0.29). Evo2 uses T15\_xs, with 547/3,053 positives (about 0.18), where deployed CNN-v2 AP is 0.726 rather than 0.848. NTv3's high historical score has unresolved overlap/provenance limits.

\begin{table}[h]
\centering
\caption{Historical curated-test average precision (higher is better within a fixed evaluation
distribution). Splits, checkpoints, input widths and weighting are not identical across rows;
this is checkpoint context, not a controlled six-model ranking.}
\label{tab:balanced}
\begin{tabular}{lc}
\toprule
Model & Historical test average precision \\
\midrule
Nucleotide Transformer v3 & 0.960 \\
ProkBERT                  & 0.884 \\
Nucleotide Transformer    & 0.854 \\
Fusion CNN                & 0.848 \\
Simple CNN                & 0.804 \\
Evo2 probe                & 0.703 \\
\bottomrule
\end{tabular}
\end{table}

\section{Standalone restricted benchmark}
Evo2 leads this report-grid comparison. Equal stride does not guarantee the same reverse-strand grid; its standalone and paired deployment values are separate runs.

\begin{table}[h]
\centering
\caption{Restricted candidate-region, step-10 benchmark precision proxy on two \emph{Clostridium}
strains, at 10\% and 50\% published-TSS recall. Higher means fewer benchmark-unmatched peaks,
not fewer proven biological false positives. \checkmark\ marks deployed ensemble members.}
\label{tab:genome}
\begin{tabular}{lccccc}
\toprule
& & \multicolumn{2}{c}{\emph{C.\ ljungdahlii}} & \multicolumn{2}{c}{\emph{C.\ autoethanogenum}} \\
\cmidrule(lr){3-4}\cmidrule(lr){5-6}
Model & Deployed & @10\% & @50\% & @10\% & @50\% \\
\midrule
Evo2 probe                & \checkmark & \textbf{71.0\%} & \textbf{58.0\%} & \textbf{85.8\%} & \textbf{71.4\%} \\
ProkBERT                  & \checkmark & 64.4\% & 47.7\% & 76.7\% & 64.6\% \\
Fusion CNN                & \checkmark & 64.4\% & 47.2\% & 81.0\% & 66.0\% \\
Simple CNN                & ---        & 49.0\% & 40.3\% & 72.6\% & 63.0\% \\
Nucleotide Transformer    & ---        & 43.9\% & 36.1\% & 64.9\% & 58.7\% \\
Nucleotide Transformer v3 & ---        & 35.7\% & 28.0\% & 58.5\% & 48.3\% \\
\bottomrule
\end{tabular}
\end{table}

\section{Paired ensemble comparison}
Deployment intersects coordinates before computing member ranks, then averages average-tie percentiles. These are not probabilities or unanimous votes; the mine uses no calibrator.

\begin{table}[h]
\centering
\caption{Three-model ensemble versus its members: restricted step-10 benchmark precision proxy
at 10\% and 50\% published-TSS recall, on the deployed shared window grid.}
\label{tab:ensemble}
\begin{tabular}{lcccc}
\toprule
& \multicolumn{2}{c}{\emph{C.\ ljungdahlii}} & \multicolumn{2}{c}{\emph{C.\ autoethanogenum}} \\
\cmidrule(lr){2-3}\cmidrule(lr){4-5}
& @10\% & @50\% & @10\% & @50\% \\
\midrule
\textbf{Three-model ensemble} & \textbf{74.4\%} & \textbf{61.6\%} & \textbf{87.0\%} & \textbf{76.6\%} \\
Evo2 probe                    & 69.4\% & 58.4\% & 87.6\% & 72.5\% \\
ProkBERT                      & 64.4\% & 47.7\% & 76.7\% & 64.6\% \\
Fusion CNN                    & 64.4\% & 47.2\% & 81.0\% & 66.0\% \\
\bottomrule
\end{tabular}
\end{table}

The trio leads at 50\% recall on both strains. Paired Evo2 leads at 10\% recall in C. autoethanogenum, 87.6\% versus 87.0\%. Adding CNN-v1 reduced the reported proxy, including 61.6\% to 58.3\% on C. ljungdahlii at 50\% recall. This supports the tested combination, not a universal optimum of three members.

\section{Deployment and limits}
We retain fine-tuned ProkBERT, from-scratch fusion CNN-v2 and the frozen-Evo2 probe. Shared training data does not establish independent errors. Deployed weights retain 888 wrong-assembly B. licheniformis positives. Cluster-disjoint splitting does not remove every dependence.

Original inputs are not committed with this report; authenticated checkpoint recovery and separate local run recovery are documented in docs/MODELS.md. New DNA scans need matching reference/grid and Evo2 runtime. The former constant-only NTv3 metric writer is not executable evaluation evidence. Ranks and motif placements prioritize experiments, not activity, strength or regulation.
\end{document}
